% Round 18. U2 branch. Candidate until independent review is bound.
\section{Inertial numerical triviality and ramified norm detection}
\label{r18ri:section}

This argument strengthens the ramified-extension conclusion by using
R17's exact necessary-and-sufficient period detector.  It needs no
component middle positivity and no toroidal weak factorization.
An arbitrary new cycle is tested by its norm, after proving that the
Galois action on every numerical model test is trivial.  That last
assertion is an actual normal-torus calculation on a Galois-equivariant
semistable resolution, not an assumption that the new cycle descends.

\subsection{A global connected-torus action on each new component}

Let $\mathcal X/R$ be a regular projective reduced strict semistable
model of a smooth projective geometrically integral $X/K$ of relative
dimension $n$.  Let $L/K$ be finite, Galois and totally ramified, with
group $G$ and ramification index $e=[L:K]$.  Write $t$ for a
uniformizer of $R_L$ and $k$ for the common residue field.  Fix an
ordering of the native components $D_1,\ldots,D_m$ of $\mathcal X_k$,
with ideals $I_i=\mathcal O_{\mathcal X}(-D_i)$.

By Saito, \emph{Weight spectral sequences and independence of $\ell$},
Lemmas 1.2.4--1.2.5, the actual blowup
\[
 \mathcal Y=
 \operatorname{Bl}_{\mathfrak a_e}(\mathcal X_{R_L}),\qquad
 \mathfrak a_e=
 \prod_{0\le j<l\le m}\ \prod_{h=1}^{e-1}
 \bigl(I_{j+1}\cdots I_l\mathcal O_{\mathcal X_{R_L}}
                        +t^h\mathcal O_{\mathcal X_{R_L}}\bigr)
 \tag{RI18.1}\label{r18ri:model}
\]
is projective reduced strict semistable over $R_L$ and has generic
fiber $X_L$.  This is the primary-source construction, valid also
for wild ramification:
\url{https://www.ms.u-tokyo.ac.jp/~t-saito/pp/wmr2r.pdf},
printed pp.~15--16.
Every factor ideal in \eqref{r18ri:model} is $G$-stable, since each
$I_i$ is defined over $R$ and $g(t)/t$ is a unit.  Hence $G$ acts
semilinearly on this actual blowup, and acts $k$-linearly on its
special fiber.

\begin{lemma}[Inertia is a connected normal-torus translation]
\label{r18ri:torus}
Every native irreducible component $F$ of $\mathcal Y_k$ is fixed
by $G$.  For each $g\in G$, the induced automorphism of $F$ is
translation by a $k$-point of a split connected algebraic torus
acting on the \emph{whole actual projective component} $F$.
Consequently
\[
                       g^*=1\quad\hbox{on }N^a(F)_{\mathbf Q}
                         \qquad\hbox{for every }a.
 \tag{RI18.2}\label{r18ri:numerical-triviality}
\]
The statement includes native components with a nontrivial finite
constant field, with their native degrees unchanged.
\end{lemma}

\begin{proof}
\emph{Step 1: the component label and its normal torus.}
A component of the Saito model is labelled by an old connected
stratum $D_I$, $s=|I|$, and a height-one ray
\[
                 ((a_i)_{i\in I},1),\qquad
                   a_i>0,\quad \sum_{i\in I}a_i=e.
 \tag{RI18.3}\label{r18ri:ray}
\]
Here $I$ is exactly the support of the ray.  The local lattice is
\[
 M_I=(\mathbf Z^I\oplus\mathbf Z[t])/
                              \mathbf Z((1)_{i\in I},-e).
\]
The characters orthogonal to the ray have rank $s-1$ and are
\emph{integrally} generated by
\[
                         z_i=x_i/t^{a_i}\quad (i\in I),
                   \qquad \prod_{i\in I}z_i=\text{a base unit}
                     \quad\hbox{over }D_I^\circ.
 \tag{RI18.4}\label{r18ri:normal-characters}
\]
Indeed the corresponding lattice generators
$[x_i]-a_i[t]$ have the single relation that their sum is zero;
any $s-1$ of them are an integral basis.
Thus the open fiber is a torsor for the split connected torus
\[
                T_I=\ker\bigl((\mathbf G_m)^I
                                         \xrightarrow{\prod}\mathbf G_m\bigr).
 \tag{RI18.5}\label{r18ri:torus-group}
\]
Its base is the actual old open stratum
$D_I^\circ=D_I\setminus\bigcup_{j\notin I}D_j$.
The new component maps to $D_I$, and its dense open normal torus
has no finite covering of that base.  This follows directly from
the preceding integral basis, not just a dimension count.

\emph{Step 2: extension of the torus action over the full boundary.}
We check extension and gluing, since translation on a dense open
alone would not suffice for \eqref{r18ri:numerical-triviality}.
Near an old stratum $D_{I\cup J}$, with $J$ the extra boundary
indices, strict semistability gives actual etale monomial charts
\[
                  \prod_{i\in I}x_i\prod_{j\in J}x_j
                                     =u\,t^e.
 \tag{RI18.6}\label{r18ri:boundary-chart}
\]
The coordinates $x_j$, $j\in J$, restrict to the old boundary
coordinates on $D_I$; the extra smooth coordinates also belong to
that base.  On the component $F$, the local orbit-character
lattice fits into the exact sequence
\[
 0\longrightarrow\mathbf Z^J
 \longrightarrow M_{I\cup J}\cap a^\perp
 \longrightarrow M_I\cap a^\perp\longrightarrow0.
 \tag{RI18.7}\label{r18ri:boundary-lattice}
\]
The left basis consists of the $[x_j]$, and the right lattice has
the integral basis just exhibited.  Choose their fixed integral
lifts.  Scaling $x_i$ for $i\in I$ by $\lambda_i$, where
$\prod_{i\in I}\lambda_i=1$, and fixing $t$, the $x_j$ and all
base functions preserves \eqref{r18ri:boundary-chart}.
Every Saito chart is a monomial subalgebra of the common character
algebra.  Such a diagonal scaling preserves that subalgebra and
all its relations, so defines an action on the entire affine chart
of $F$, including its toric boundary.  This is also immediate in
the subdivision description: a torus preserves every toric orbit
chart of its fan.  Only the normal characters are scaled; the old
boundary and smooth base functions are fixed.

These chart actions glue on the \emph{actual} component.  On an
overlap, a local equation of $D_i$ changes from $x_i$ to $v_i x_i$
with $v_i$ invertible.  Since the map $F\to\mathcal X_k$ factors
through $D_I$, its restriction to $F$ is the pullback of
$v_i|_{D_I}$.  Thus normal coordinates change by multiplication by
units from the base.  These transitions commute with the same
$T_I$ action.  Transitions involving old boundary coordinates are
integral monomials with coefficients in base units and preserve
the character-lattice description \eqref{r18ri:boundary-lattice}.
Here is the precise etale descent check.  Let $U\to U_0$ be one
of the actual etale charts of the old model, with $U_0$ the standard
monomial scheme, and let $C_0$ be an affine Saito chart over
$(U_0)_{R_L}$.  Blowup commutes with this flat base change.
If $F_0\subset C_0$ is the ray component, its structural map
factors scheme-theoretically through the old coordinate stratum
$(D_I)_0$, because every $x_i$ with $i\in I$ has positive valuation
and vanishes on the reduced $F_0$.  Therefore the corresponding
component chart over $U$ is
\[
 F_0\times_{(U_0)_k}U_k
       =F_0\times_{(D_I)_0}(D_I\cap U).
\]
The diagonal $T_I$ action on $F_0$ fixes $(D_I)_0$, so it acts on
this actual fiber product.  This is a base change over the old
stratum, not an assertion that an arbitrary etale covering of a
torus-equivariant scheme inherits its action.  On overlaps the
homogeneous monomial generators change by the base units already
computed, while every old regular function restricts through
$D_I$.  Thus their coactions agree.  This also proves the assertion
in a neighborhood of the entire old boundary, rather than only
over $D_I^\circ$.  No formal normal neighborhood is substituted
for an actual algebraic chart.  This proves a global algebraic
$T_I$ action on $F$.

\emph{Step 3: the inertia translation and numerical action.}
For $g\in G$, put
\[
                              c_g=\overline{g(t)/t}\in k^\times.
\]
Total ramification makes $g$ the identity on the residue field and
on every old stratum.  From $\pi=u_0t^e$ one gets $c_g^e=1$:
$g(u_0)/u_0$ has residue one.  Every local $x_i$ is defined on the
original model and is fixed by $g$.  Thus on the dense normal
torus of $F$,
\[
                         g(z_i)=c_g^{-a_i}z_i.
 \tag{RI18.8}\label{r18ri:translation}
\]
The scalars have product $c_g^{-e}=1$, so they give a point of
$T_I(k)$.  The ray and old connected stratum labels are fixed, so
$F$ is fixed.  The two automorphisms in question, the given inertia
action and this global torus translation, agree on a dense open
of the integral separated variety $F$; hence they agree everywhere.
In particular the part with $c_g=1$, including wild inertia, acts
identically on the component.

Finally, translation by a point of a split connected torus acts
trivially on numerical cycle groups of a projective variety.
One can verify this without a cycle-class comparison: move the
point one coordinate at a time along $\mathbf G_m$.  For any
cycle, its translates form a proper flat family of cycles over
that connected curve, obtained from the product family by the
action automorphism.  Intersection degrees against every fixed
complementary class are constant.  The endpoints are therefore
numerically equivalent.  This proves
\eqref{r18ri:numerical-triviality}.  All these are native degrees on
proper $k$-varieties; a finite constant field only multiplies them
by its full positive degree.
\end{proof}

\subsection{Norm detection by the exact model-period kernel}

Fix a codimension $1\le r\le n$ and put $a=n-r$.
On the native reduced strict semistable model $\mathcal Y/R_L$ use
\[
 C^a=\bigoplus_FN^a(F)_{\mathbf Q},\qquad V_a=\ker A_a,
\]
and let $\nu_Z\in V_a^\vee$ be the frozen FD17 model-period
functional of a rational cycle class $Z$ on $X_L$.
FD17 and VT17 identify its zero kernel with strong G for that
specified cycle; ER17's edge and compatible quotients do not change
that kernel.  This includes the necessary direction, not just a
sufficient numerical criterion.

\begin{theorem}[Ramified Galois norm detects the full obstruction]
\label{r18ri:norm}
In the situation above, write $p:X_L\to X$ for the finite field
projection.  For any specified rational cycle $Z$ on $X_L$,
\[
                       G_K(p_*Z)\quad\Longrightarrow\quad G_L(Z).
 \tag{RI18.9}\label{r18ri:norm-detection}
\]
In particular, if every geometrically homologically trivial
codimension-$r$ rational cycle on $X$ has strong G, then the same
is true for every such cycle on $X_L$, including cycles which do
not descend to $K$.
\end{theorem}

\begin{proof}
The Galois action preserves the actual model, its component labels,
and all native intersection operations.  For every $v\in V_a$,
\[
                          \nu_{gZ}(v)=\nu_Z(g^{-1}v)=\nu_Z(v).
 \tag{RI18.10}\label{r18ri:period-equivariance}
\]
The first equality can be checked directly with an actual pure
operational Chow extension $\Theta$ of $Z$: on each component it
is the native degree projection formula for the isomorphism $g$,
and its normalized uniformizer value is unchanged because
$|g(t)|=|t|$.  The second equality is
Lemma~\ref{r18ri:torus}.  Numerical representatives suffice only
for this native intersection pairing; no unproved geometric
cycle-class map on $N^a(F)$ is used.  No degree is taken on the
total model's absolute Chow zero-cycle group.

Suppose $G_K(p_*Z)$ holds and write $B=p_*Z$ for its actual
collected rational cycle.  We spell out the finite-field ascent on
all tests.  If $\gamma_B$ is its supplied strong Green form and
$\eta$ is any compact $L$-actual test, FG16's finite-constant
comparison reads $\eta$ in the native $K$ theory on $X_L$.
Full finite actual trace and FG16's arbitrary-cycle restriction
identity give
\[
\begin{aligned}
 \int_{(X_L\setminus|p^*B|)/L}p^*\gamma_B\wedge D\eta
 &=\frac1{[L:K]}
     \int_{(X\setminus|B|)/K}
            \gamma_B\wedge D\operatorname{Tr}_p^K\eta\\
 &=-\frac1{[L:K]}\int_{B/K}\operatorname{Tr}_p^K\eta
   =-\int_{p^*B/L}\eta .
\end{aligned}
\]
The polar cutoff and exact-support argument are precisely FG16.3;
the trace is ordinary actual on every test, and the common native
integration multiplier is the full $[L:K]$.  Thus $G_L(p^*p_*Z)$
holds without any restriction to pullback tests.  Alternatively it
follows from the already frozen obstruction pullback.  Finite flat
Galois base change gives the actual Chow identity
\[
                             p^*p_*Z=\sum_{g\in G}gZ.
 \tag{RI18.11}\label{r18ri:norm-identity}
\]
By the necessary direction of the exact FD17/VT17 detector, the
period of this cycle vanishes on every $v\in V_a$.
Using \eqref{r18ri:period-equivariance} and linearity gives
\[
                   0=\nu_{p^*p_*Z}(v)=|G|\,\nu_Z(v).
 \tag{RI18.12}\label{r18ri:averaging}
\]
The full nonzero field degree $|G|=[L:K]$ is invertible with
rational coefficients, including its wild part.  Hence
$\nu_Z=0$.  The sufficient direction of the same exact detector
and the frozen reconstruction give the Green form for the
\emph{original specified} cycle $Z$ on its exact support complement,
with ordinary actual output and every original compact actual test.
No assertion that $v$ represents every compact test is needed:
its vanishing is only the obstruction test; the final reconstructed
identity retains the full test space.

For the last assertion, finite Chow pushforward takes a geometrically
homologically trivial cycle to a geometrically homologically trivial
cycle: after passing to the geometric generic fiber its class is the
sum of the conjugate zero classes.  Thus the hypothesis on $X$
applies to the norm of each arbitrary new cycle over $L$.
\end{proof}

\begin{corollary}[From all unramified fields to all finite fields]
\label{r18ri:all-finite}
Let $X/K$ have the given native reduced strict semistable projective
model.  Suppose that, for every finite unramified extension $K'/K$,
every geometrically homologically trivial rational cycle of a fixed
codimension $r$ on $X_{K'}$ has strong G.  Then the same full
codimension-$r$ assertion holds after every finite valued extension
$L/K$.
\end{corollary}

\begin{proof}
Choose a finite Galois closure $E/K$ containing $L$, and let
$K_0=E^{I}$ be the maximal unramified subextension, where $I$ is
the inertia subgroup.  Then $E/K_0$ is totally ramified and Galois.
The model over $K_0$ is still native reduced strict semistable,
by unramified base change.  The assumed full theorem over $K_0$
and Theorem~\ref{r18ri:norm} give the full theorem over $E$.
Given an arbitrary specified cycle over $L$, pull it to $E$,
apply that theorem, and descend by the frozen finite actual trace
with the complete degree $[E:L]$ and the original absolute value
on $L$.  This returns the final exact-support all-actual Green
form over $L$.  No cycle over $L$ was assumed to descend to $K_0$
or to $K$.
\end{proof}

\subsection{Consequences for the new product families}

\begin{corollary}[Every finite extension for the complete product-model domains]
\label{r18ri:product-families}
For every $n\ge4$, every native nonlinear-factor or hyperplane
model covered by GH18 has full strong G in every codimension after
\emph{every} finite valued extension, for every new specified cycle.
For the GF18 and NG18 dense global families, the complete
complementary-degree height pairings of GH18 therefore exist over
\emph{every} number field $F$, with no unramified-at-$q$ restriction:
\[
 \operatorname{CH}^r(X_F)_{\mathbf Q,\mathrm{geom\,hom}}
 \ \times\
 \operatorname{CH}^{n+1-r}(X_F)_{\mathbf Q,\mathrm{geom\,hom}}
                                   \longrightarrow\mathbf R
                 \qquad(1\le r\le n).
 \tag{RI18.13}\label{r18ri:all-heights}
\]
\end{corollary}

\begin{proof}
GH18 supplies full G on every unramified extension by applying its
model theorem anew to every cycle over that field.  The model is
already reduced strict semistable, so
Corollary~\ref{r18ri:all-finite} applies in each codimension.
In the global families, apply this result at every place above the
distinguished $q$; at every other place the frozen good or
ordinary-quadratic theorem already accepts arbitrary finite
extension.  The full local hypotheses of PH16 hold for every
complementary pair over $F$, and GA17 supplies admissible lifts on
this established domain.  This proves \eqref{r18ri:all-heights}.
The coefficient families, their density and their monodromy ranks
are unchanged; no new spreading or density assertion is required.
\end{proof}

\paragraph{Scope and dependencies.}
This theorem propagates an already proved full cycle domain from all
unramified extensions to all finite extensions of a variety with a
supplied strict semistable model.  It does not prove the starting
vanishing for an arbitrary variety, does not prove that a sufficient
model quotient must vanish, and does not settle general G or arbitrary
full BB heights.  Its indispensable inputs are the actual
Galois-equivariant Saito blowup, the global normal-torus action proved
above, the exact necessary direction of FD17/VT17, ordinary Chow norm
and pullback, and the frozen finite actual pullback/trace and
exact-support reconstruction.  There is no arbitrary proper actual
pushforward, no component Hodge conjecture, and no use of the separate
RB18 weak-factorization argument.  No literature-priority claim is made.
